\section*{Author Section 2.3: quantum communication and games}
A quantum protocol for $f\subseteq\mathcal X\times\mathcal Y\times\mathcal Z$ has Alice and Bob with inputs $x,y$ and prior shared entanglement. Each round one party applies a unitary to its classical input, last received message and retained memory, keeps some registers and sends the rest. Safe copies of inputs can be made by CNOT so input registers emerge unchanged. Shared classical randomness is included in this model by both parties measuring shared entanglement in the same basis. In one-way communication Alice sends the single message; retained registers are discarded. Bob performs a projective measurement on his input, shared entanglement and message, yielding the output $P(x,y)$. Equivalently he performs a unitary then measures designated $\log|\mathcal Z|$ output qubits in the computational basis. Arbitrary entanglement and local-operation complexity are allowed; the resource is qubits sent.
\paragraph{Definition 26.} For $0<\epsilon<1$, $Q^1_\epsilon(f)$ is the minimum message size of a one-way entanglement-assisted protocol with $\Pr[P(x,y)\in f(x,y)]\ge1-\epsilon$ for every $(x,y)$, over inherent protocol randomness.
\paragraph{Definition 27.} $Q^1_{p,\epsilon}(f)$ is the corresponding minimum message size when success $\ge1-\epsilon$ is averaged over $(x,y)\sim p$ and protocol randomness.
\paragraph{Fact 28 (Yao's lemma [45]).} $Q^1_\epsilon(f)=\max_pQ^1_{p,\epsilon}(f)$ for $0<\epsilon<1$.
A two-player non-local game is $G=(q,\mathcal X\times\mathcal Y,\mathcal A\times\mathcal B,V)$: the referee samples $(x,y)\sim q$, gives inputs, and accepts outputs iff $V(x,y,a,b)=1$. Entangled players measure their respective shared-state halves and inputs (unitary followed by computational-basis measurement on designated $\log|\mathcal A|$, $\log|\mathcal B|$ output qubits).
\paragraph{Definition 29.} The entangled value $\omega^*(G)$ is the maximum winning probability over all shared-entanglement strategies, averaged over $q$ and inherent randomness. A distribution is anchored on one side with probability $\zeta$ if either some $x^*$ has $q(x^*)=\zeta$, $q(y|x^*)=q(y)$ for all $y$, or some $y^*$ has $q(y^*)=\zeta$, $q(x|y^*)=q(x)$ for all $x$ (author Section1.1).
